% ===================== Chapter 13: NumPy (ARC500) =====================
% Requires your project preamble to define:
% \h{...}, \hh{...}, \hhh{...}
% \code{<label>}{python}{<code>}{<caption>}
% All Python examples below use EXACTLY 4 spaces per indent (no tabs).

\h{NumPy (Numerical Python)}

\emph{NumPy} is the foundation for fast, array-based computing in Python. It stores numbers in compact, contiguous blocks and applies math to \emph{whole arrays} at once (vectorization). For architects and structural designers, this maps naturally to grids, room schedules, facade panels, load combinations, and environmental datasets.

In practice you will:
\begin{itemize}
\item build and reshape 1D/2D arrays representing spans, panel centers, or floor areas;
\item index/slice subsets (levels, bays, wings) and filter with boolean conditions;
\item broadcast unit conversions and design formulas across whole arrays;
\item aggregate by axes (sum by floor, average by bay), sort/search, join/split datasets;
\item apply universal functions (ufuncs) like \c{add}, \c{log}, \c{diff}, \c{lcm}, \c{sin}/\c{cos}.
\end{itemize}

\begin{tcolorbox}[% passing options here
    title=Focus of this chapter:,
    colframe=orange,
    colback = white,
%    colback=orange!40!yellow!30, % think of mixing 2 colors with sliders
    arc=2mm % sharp corners
 ]
\begin{itemize}
\item create arrays (\c{array}, \c{arange}, \c{linspace}, \c{zeros}, \c{ones}, \c{full}, \c{eye});
\item read basics: dtype, shape, ndim;
\item index/slice and filter with boolean masks and \c{np.isin};
\item distinguish view vs copy and avoid accidental coupling;
\item reshape, transpose, flatten; use broadcasting correctly;
\item compute aggregates by axis; sort, search, join, split;
\item use core ufuncs (arithmetic, logs, products/sums, \c{diff}, \c{lcm}, trig).
\end{itemize}
\end{tcolorbox}

% ====================================================================
% PART I — NUMPY BASICS
% ====================================================================

\hh{Getting started (install and import)}

NumPy lives outside the Python standard library. You install it once (on your machine), and then import it at the top of any script that uses arrays. The alias \c{np} is the community standard—you will see it everywhere.

\code{c14_setup}{python}{
    # In a terminal (not inside Python):
    # pip install numpy

    import numpy as np
    print(np.__version__)
    # Outcome:
    # e.g., 2.0.x
}{Check your NumPy version}

\hh{Arrays vs lists (why NumPy is fast)}

Python lists can hold anything (numbers, strings, even mixed types), but that flexibility makes them slow for math. NumPy arrays are homogeneous numeric blocks; the library uses optimized C loops to apply one operation to all elements at once. Think “operate on the whole schedule, not one cell at a time.”

\code{c14_list_vs_array}{python}{
    import numpy as np

    widths  = [3.5, 4.0, 5.0]     # Python list (m)
    lengths = [8.0, 7.5, 6.0]     # Python list (m)

    w = np.array(widths)          # NumPy arrays
    L = np.array(lengths)
    areas = w * L                 # elementwise multiply (vectorized)
    print(areas)
    # Outcome:
    # [28. 30. 30.]
}{Elementwise math with arrays}

\hh{Creating arrays}

There are many ways to make arrays. Use \c{array([...])} to convert existing Python data; \c{arange} or \c{linspace} to create number sequences; and “fillers” like \c{zeros}/\c{ones}/\c{full}/\c{eye} for standard shapes. Choosing the right constructor keeps code short and clear.

\code{c14_make_arrays}{python}{
    import numpy as np
    np.random.seed(0)  # reproducible example

    a = np.array([2, 4, 6])               # from a Python list
    b = np.arange(0, 10, 2)               # 0,2,4,6,8
    c = np.linspace(0, 1, 5)              # 5 points from 0..1
    Z = np.zeros((2, 3))                  # 2x3 zeros (float)
    O = np.ones((3, 2), dtype=int)        # 3x2 ones (int)
    F = np.full((2, 2), 3.14)             # 2x2 filled with 3.14
    I = np.eye(3)                         # 3x3 identity
    R = np.random.rand(2, 3)              # random [0,1)

    print(a, b, c, sep="\\n")
    print(Z, O, F, sep="\\n")
    print(I, R, sep="\\n")
    # Outcome:
    # Arrays with the shapes/types as created
}{From lists and helper constructors}

\hh{Array basics: dtype, shape, ndim}

Every array has a \emph{shape} (rows, columns), a number of dimensions (\c{ndim}), and a data type (\c{dtype}). Read these first—like checking units and grid size before you calculate. Misunderstood shapes cause most beginner bugs.

\code{c14_basics}{python}{
    import numpy as np

    grid = np.array([[6.0, 6.0, 6.0],
                     [7.5, 7.5, 7.5]])
    print("shape:", grid.shape)  # rows x cols
    print("ndim:", grid.ndim)    # number of axes
    print("dtype:", grid.dtype)  # element type
    # Outcome:
    # shape: (2, 3)
    # ndim: 2
    # dtype: float64
}{Read array metadata quickly}

\hh{Indexing and slicing}

Indexing lets you read or write individual cells; slicing selects ranges or entire rows/columns. Negative indices count from the end; the slice form \c{start:stop:step} behaves like CAD ranges (end not included). Practice 1D and 2D patterns—they are the foundation for everything else.

\code{c14_index_slice}{python}{
    import numpy as np

    bays = np.array([6.0, 7.5, 9.0, 7.5, 6.0])  # bay spans (m)
    print(bays[0])        # first
    print(bays[-1])       # last
    print(bays[1:4])      # slice indices 1..3
    print(bays[::2])      # every second

    floors = np.array([[120, 130, 125],
                       [140, 150, 145]])        # areas by wing (m^2)
    print(floors[0, 1])   # row 0, col 1
    print(floors[:, 0])   # first column
    print(floors[1, :])   # second row
    # Outcome:
    # 6.0
    # 6.0
    # [7.5 9.  7.5]
    # [6. 9. 6.]
    # 130
    # [120 140]
    # [140 150 145]
}{1D/2D indexing patterns}

\hh{Filtering with boolean masks (where, isin)}

Masks are boolean arrays (True/False) that pick elements matching a condition—like querying a schedule (“area >= 30”). Combine conditions with \c{(cond1) \& (cond2)} and use \c{np.isin} to select from a set. Remember: \c{\&}, \c{|}, \c{~} (bitwise) replace Python’s \c{and}, \c{or}, \c{not} for elementwise logic.

\code{c14_boolean_masks}{python}{
    import numpy as np

    areas_m2 = np.array([28.0, 32.5, 18.0, 45.0, 30.0])
    large = areas_m2 >= 30.0
    print("mask:", large)
    print("large rooms:", areas_m2[large])

    bays = np.array([6.0, 7.5, 9.0, 7.5, 6.0])
    keep_values = np.array([7.5, 9.0])
    mask = np.isin(bays, keep_values)
    print("keep mask:", mask)
    print("selected bays:", bays[mask])
    # Outcome:
    # mask: [False  True False  True  True]
    # large rooms: [32.5 45.  30. ]
    # keep mask: [False  True  True  True False]
    # selected bays: [7.5 9.  7.5]
}{Select by condition or set}

\hh{Copy vs view (avoid accidental coupling)}

Slices usually return \emph{views} (new array object, same data). Editing a view edits its parent—great for performance, surprising when you expect independence. Call \c{.copy()} to break the link. A good habit: copy results you plan to modify.

\code{c14_copy_vs_view}{python}{
    import numpy as np

    base = np.array([10.0, 20.0, 30.0, 40.0])
    window = base[1:3]        # view
    window[:] = 99.0          # modifies base too
    print("base after editing view:", base)

    base2 = np.array([10.0, 20.0, 30.0, 40.0])
    safe = base2[1:3].copy()  # independent copy
    safe[:] = 77.0
    print("base2:", base2)
    print("safe copy:", safe)
    # Outcome:
    # base after editing view: [10. 99. 99. 40.]
    # base2: [10. 20. 30. 40.]
    # safe copy: [77. 77.]
}{Slicing returns a view; use {copy()} to decouple}

\hh{Shape, reshape, transpose, flatten}

Reshaping and transposing are “different views of the same data.” They let you reinterpret layout without duplicating memory. Use \c{ravel()} to flatten as a view when possible; prefer it over \c{flatten()} (which always copies) unless you truly need a new buffer.

\code{c14_reshape}{python}{
    import numpy as np

    rooms = np.arange(12)          # 0..11 (room ids)
    plan = rooms.reshape(3, 4)     # 3 rows, 4 cols
    print(plan)

    flat = plan.ravel()            # view if possible
    transposed = plan.T            # 4x3
    print(flat)
    print(transposed)
    # Outcome:
    # [[ 0  1  2  3]
    #  [ 4  5  6  7]
    #  [ 8  9 10 11]]
    # [ 0  1  2  3  4  5  6  7  8  9 10 11]
    # [[ 0  4  8]
    #  [ 1  5  9]
    #  [ 2  6 10]
    #  [ 3  7 11]]
}{Reshape views of the same data}

\hh{Broadcasting (unit-wise math without loops)}

Broadcasting lets arrays with compatible shapes interact without looping or copying. Dimensions are compatible when they are equal or one of them is 1. Picture a column stretched across rows or a row stretched across columns—like applying a single conversion factor to an entire facade grid.

\code{c14_broadcasting}{python}{
    import numpy as np

    widths  = np.array([3.0, 3.5, 4.0])  # (3,)
    lengths = np.array([[8.0], [9.0]])   # (2,1)
    areas = lengths * widths             # -> (2,3)
    print(areas)

    m2_to_ft2 = 10.7639
    areas_ft2 = areas * m2_to_ft2
    print(areas_ft2)
    # Outcome:
    # [[24.  28.  32. ]
    #  [27.  31.5 36. ]]
    # [[258.3336 301.3892 344.4448]
    #  [290.6453 339.4759 387.3064]]
}{Shape compatibility for whole-array math}

\hh{Vectorized arithmetic and aggregation}

Vectorized formulas read like the math you would write on paper. Aggregations summarize by axis: \c{axis=0} collapses rows (column-wise totals); \c{axis=1} collapses columns (row-wise totals). Always track units (m, m\textsuperscript{2}, kN/m\textsuperscript{2}) in your variable names or comments.

\code{c14_arith_agg}{python}{
    import numpy as np

    w = np.array([3.0, 3.5, 4.0])
    L = np.array([8.0, 7.5, 6.0])

    A = w * L                      # areas
    P = 2 * (w + L)                # perimeters
    print("A:", A)
    print("P:", P)

    print("total area:", A.sum())
    print("mean width:", w.mean())
    print("max length:", L.max())
    # Outcome:
    # A: [24.   26.25 24.  ]
    # P: [22. 22. 20.]
    # total area: 74.25
    # mean width: 3.5
    # max length: 8.0
}{Express design math clearly}

\hh{Sorting and searching}

Sorting reorders values; \c{argsort} returns the indices that would sort (useful for reordering multiple aligned arrays). \c{where} finds positions meeting a condition. \c{searchsorted} inserts into a sorted array—handy for locating gridline positions for new elements.

\code{c14_sort_search}{python}{
    import numpy as np

    A = np.array([32.5, 18.0, 45.0, 30.0])
    print("sorted:", np.sort(A))

    idx = np.argsort(A)
    print("argsort indices:", idx)
    print("use indices:", A[idx])

    print("where >= 30:", np.where(A >= 30.0))

    gridlines = np.array([0, 3, 6, 9, 12])   # meters
    print("insert 7m at index:", np.searchsorted(gridlines, 7))
    # Outcome:
    # sorted: [18.  30.  32.5 45. ]
    # argsort indices: [1 3 0 2]
    # use indices: [18.  30.  32.5 45. ]
    # where >= 30: (array([0, 2, 3]),)
    # insert 7m at index: 3
}{Order and locate values}

\hh{Joining and splitting arrays}

“Joining” stacks arrays together; “splitting” cuts one into equal parts. Use \c{concatenate} along the matching axis (shapes must align). These are the basic mechanics for assembling multi-wing datasets or breaking a long vector into floors.

\code{c14_join_split}{python}{
    import numpy as np

    wing_A = np.array([120, 130, 125])
    wing_B = np.array([110, 140, 135])

    whole = np.concatenate([wing_A, wing_B])
    print("joined:", whole)

    floors = np.arange(12)   # 0..11
    parts = np.split(floors, 3)
    print("split parts:", parts)
    # Outcome:
    # joined: [120 130 125 110 140 135]
    # split parts: [array([0, 1, 2, 3]), array([4, 5, 6, 7]), array([ 8,  9, 10, 11])]
}{Glue datasets or break evenly}

\hh{Iterating (only when necessary)}

Prefer array math to Python \c{for} loops: it is faster, clearer, and less error-prone. If you must loop (e.g., I/O per element), \c{np.nditer} gives a consistent iterator over any shape; still try to keep heavy computation vectorized.

\code{c14_iterating}{python}{
    import numpy as np

    plan = np.array([[28.0, 30.0, 26.0],
                     [32.5, 29.0, 31.0]])

    for x in np.nditer(plan):
        print(float(x))
    # Outcome:
    # 28.0
    # 30.0
    # 26.0
    # 32.5
    # 29.0
    # 31.0
}{Iteration is slower than vectorization}

\hh{Data types (dtype) and casting}

Dtypes control storage and math behavior. Integers truncate, floats carry decimals, and mixing types promotes to a common type. Use \c{astype} to be explicit (e.g., convert to float before dividing). Be mindful of integer division if you expect decimals.

\code{c14_dtype_cast}{python}{
    import numpy as np

    x = np.array([1, 2, 3])                  # int
    y = x.astype(float)                      # to float
    z = np.array([1.2, 2.5, 3.8], dtype=int) # truncates toward 0
    print(x.dtype, y.dtype, z)
    # Outcome (platform-dependent names):
    # int64 float64 [1 2 3]
}{Be explicit when needed}

\hh{Array shape and reshape — more patterns}

Use \c{-1} to ask NumPy to infer one dimension automatically (based on the total number of elements). Reshapes must preserve the element count; if they do not, NumPy raises an error—just like trying to fit 13 panels into a 3x4 grid.

\code{c14_shape_more}{python}{
    import numpy as np

    bays = np.arange(1, 13)            # 1..12
    grid3x4 = bays.reshape(3, 4)
    grid4x3 = bays.reshape(4, 3)
    auto    = bays.reshape(3, -1)      # infer second dim
    print(grid3x4.shape, grid4x3.shape, auto.shape)
    # Outcome:
    # (3, 4) (4, 3) (3, 4)
}{Use -1 to infer a dimension}

% ====================================================================
% PART II — UFUNCS (UNIVERSAL FUNCTIONS)
% ====================================================================

\hh{Ufuncs: fast elementwise math}

Ufuncs are NumPy’s “universal functions.” They compute elementwise across arrays, follow broadcasting rules, and are heavily optimized. Prefer ufuncs over manual loops; they are both faster and clearer.

\hhh{Simple arithmetic}

These four are the bread and butter; they mirror \c{+}, \c{-}, \c{*}, \c{/} but work cleanly with broadcasting and optional arguments like \c{out=}.

\code{c14_ufunc_basic}{python}{
    import numpy as np

    x = np.array([1.0, 2.0, 3.0])
    y = np.array([10.0, 20.0, 30.0])

    print(np.add(x, y))       # x + y
    print(np.subtract(y, x))  # y - x
    print(np.multiply(x, y))  # x * y
    print(np.divide(y, x))    # y / x
    # Outcome:
    # [11. 22. 33.]
    # [ 9. 18. 27.]
    # [10. 40. 90.]
    # [10. 10. 10.]
}{Elementwise arithmetic ufuncs}

\hhh{Logs (scales, e.g., acoustics)}

Logarithms convert multiplicative scales to additive ones. Use \c{log} (natural) and \c{log10} (base-10). For plotting performance curves or sound levels, logs linearize your data.

\code{c14_ufunc_logs}{python}{
    import numpy as np

    levels = np.array([1, 10, 100, 1000], dtype=float)
    print(np.log(levels))    # natural log
    print(np.log10(levels))  # base-10 log
    # Outcome:
    # [0.         2.30258509 4.60517019 6.90775528]
    # [0. 1. 2. 3.]
}{Logarithms for scaling}

\hhh{Products and summations (by axis)}

Axis arguments control “which way you fold the array.” Summing by floor vs by bay is just \c{axis=1} vs \c{axis=0}. Learn this once; use it everywhere.

\code{c14_ufunc_prod_sum}{python}{
    import numpy as np

    loads = np.array([[3.0, 3.2, 3.1],  # kN/m^2
                      [2.8, 3.0, 3.3]])

    print("sum all:", np.sum(loads))
    print("sum by floor (axis=1):", loads.sum(axis=1))
    print("sum by bay   (axis=0):", loads.sum(axis=0))
    print("prod row 0:", loads[0].prod())
    # Outcome:
    # sum all: 18.4
    # sum by floor (axis=1): [9.3 9.1]
    # sum by bay   (axis=0): [5.8 6.2 6.4]
    # prod row 0: 30.816000000000003
}{Aggregate along axes}

\hhh{Differences (spacings/gradients)}

\c{np.diff} computes adjacent differences—perfect for story heights from floor elevations, panel spacing from gridlines, or gradient checks.

\code{c14_ufunc_diff}{python}{
    import numpy as np

    elevations = np.array([0.0, 3.2, 6.5, 10.1])  # floor elevations (m)
    story_heights = np.diff(elevations)
    print(story_heights)
    # Outcome:
    # [3.2 3.3 3.6]
}{Adjacent differences}

\hhh{Least common multiple (module coordination)}

LCM helps coordinate repeating modules—like aligning 600\,mm and 750\,mm grids for a facade pattern.

\code{c14_ufunc_lcm}{python}{
    import numpy as np

    # Find a repeating module compatible with 600 mm and 750 mm grids
    print(np.lcm(600, 750), "mm")
    # Outcome:
    # 3000 mm
}{Find a shared module}

\hhh{Trigonometric (angles/orientation)}

Angles are usually in degrees in architectural documents; convert to radians for trig functions. Great for sun angles, ramp slopes, and orientation vectors.

\code{c14_ufunc_trig}{python}{
    import numpy as np

    degrees = np.array([0, 30, 45, 60, 90], dtype=float)
    radians = np.deg2rad(degrees)
    print("sin:", np.sin(radians))
    print("cos:", np.cos(radians))
    # Outcome:
    # sin: [0.         0.5        0.70710678 0.8660254  1.        ]
    # cos: [1.         0.8660254  0.70710678 0.5        0.        ]
}{Angles in degrees to sines/cosines}

% ====================================================================
% PART III — MINI-LABS (ARCHITECTURE CONTEXT)
% ====================================================================

\hh{Mini-Lab 1 — Studio room areas and filtering}

\textit{What you practice:} vectorized area computations, unit conversion, boolean masks.  
\textit{Why it matters:} turn a spreadsheet of rooms into quick metrics without loops.

\code{c14_minilab1}{python}{
    import numpy as np

    widths  = np.array([3.2, 3.4, 4.0, 3.0])   # m
    lengths = np.array([8.0, 8.5, 8.5, 8.5])   # m
    areas_m2 = widths * lengths
    areas_ft2 = areas_m2 * 10.7639
    big_mask = areas_ft2 > 350.0
    print("areas m2:", areas_m2)
    print("areas ft2:", areas_ft2)
    print("big rooms (ft2>350):", areas_ft2[big_mask])
    # Outcome (rounded):
    # areas m2: [25.6 28.9 34.  25.5]
    # areas ft2: [275.6000 311.016  365.9726 274.4925]
    # big rooms (ft2>350): [365.9726]
}{Filter by threshold}

\hh{Mini-Lab 2 — Facade panelization grid}

\textit{What you practice:} \c{arange}, \c{meshgrid}, stacking, and counting points.  
\textit{Why it matters:} generate panel centers for downstream geometry or analysis.

\code{c14_minilab2}{python}{
    import numpy as np

    W, H = 60.0, 24.0
    sx, sy = 1.2, 0.8

    xs = np.arange(sx/2, W, sx)    # center-to-center starts at half-step
    ys = np.arange(sy/2, H, sy)
    X, Y = np.meshgrid(xs, ys)     # shapes: (len(ys), len(xs))

    centers = np.column_stack((X.ravel(), Y.ravel()))
    print("panel count:", centers.shape[0])
    print("first five centers:\\n", centers[:5])
    # Outcome (depends on spacings):
    # panel count: 1500
    # first five centers:
    # [[0.6 0.4]
    #  [1.8 0.4]
    #  [3.  0.4]
    #  [4.2 0.4]
    #  [5.4 0.4]]
}{Use {meshgrid} and {column\_stack}}

\hh{Mini-Lab 3 — Load case combination}

\textit{What you practice:} broadcasting and axis-wise totals.  
\textit{Why it matters:} apply code combinations across many bays/floors at once.

\code{c14_minilab3}{python}{
    import numpy as np

    DL = np.array([[4.0, 4.2, 4.1],    # kN/m^2
                   [3.9, 4.0, 4.3]])
    LL = np.array([[2.0, 1.8, 2.2],
                   [2.1, 2.4, 2.0]])

    combo = 1.2 * DL + 1.6 * LL
    print("combo per bay:\\n", combo)
    print("floor totals:", combo.sum(axis=1))
    print("building total:", combo.sum())
    # Outcome:
    # combo per bay:
    # [[7.2  7.2  7.54]
    #  [7.02 7.44 7.48]]
    # floor totals: [21.94 21.94]
    # building total: 43.88
}{Vectorized load combinations}

\hh{Mini-Lab 4 — Copy vs view bug hunt}

\textit{What you practice:} detecting view coupling and breaking it with \c{copy()}.  
\textit{Why it matters:} prevents accidental edits to source arrays in larger pipelines.

\code{c14_minilab4}{python}{
    import numpy as np

    plan = np.array([[28.0, 30.0, 26.0],
                     [32.5, 29.0, 31.0],
                     [27.0, 33.0, 25.0]])

    wing = plan[:, 1]        # likely a view
    wing[:] = 99.0
    print("plan after edit:\\n", plan)

    # fix
    plan2 = plan.copy()
    wing_safe = plan2[:, 1].copy()
    wing_safe[:] = 77.0
    print("plan2 (safe):\\n", plan2)
    # Outcome:
    # plan after edit:
    # [[28. 99. 26.]
    #  [32.5 99. 31.]
    #  [27. 99. 25.]]
    # plan2 (safe):
    # [[28. 99. 26.]
    #  [32.5 77. 31.]
    #  [27. 99. 25.]]
}{Demonstrate view-coupling and fix with copy}

% ====================================================================
% PART IV — QUICK REFERENCE & TIPS
% ====================================================================

\hh{Key syntax (quick reference)}
\begin{itemize}
\item Create: \c{np.array([...])}, \c{np.arange(start, stop, step)}, \c{np.linspace(a, b, n)}
\item Fillers: \c{np.zeros((r,c))}, \c{np.ones((r,c), dtype=...)}, \c{np.full((r,c), value)}, \c{np.eye(n)}
\item Introspection: \c{arr.shape}, \c{arr.ndim}, \c{arr.dtype}
\item Index/slice: \c{arr[i]}, \c{arr[i:j:k]}, \c{arr[r, c]}, \c{arr[:, 0]}, \c{arr[1, :]}
\item Filter: \c{mask = arr > 0}, \c{arr[mask]}, \c{np.where(mask)}, \c{np.isin(arr, values)}
\item Views/copies: slices $\to$ view; \c{arr.copy()} for independence; \c{ravel()} view if possible; \c{flatten()} copy
\item Reshape/transpose: \c{arr.reshape(r,c)}, \c{arr.T}
\item Aggregation: \c{arr.sum()}, \c{arr.mean()}, \c{arr.sum(axis=0/1)}
\item Sort/search: \c{np.sort(x)}, \c{np.argsort(x)}, \c{np.where(cond)}, \c{np.searchsorted(sorted, v)}
\end{itemize}

\hh{Common pitfalls and tips}
\begin{itemize}
\item \textbf{Views vs copies:} Slices are often views. Use \c{.copy()} for independence. \c{ravel()} is a view if possible; \c{flatten()} always copies.
\item \textbf{Shapes must align:} Broadcasting needs size 1 or equal sizes along each dimension.
\item \textbf{Axis matters:} \c{sum(axis=0)} aggregates rows (column-wise); \c{sum(axis=1)} aggregates columns (row-wise).
\item \textbf{Prefer vectorization:} Use array math over Python \c{for} loops for speed and clarity.
\item \textbf{Be explicit with dtypes:} Mixed integer/float math promotes to float; specify \c{dtype=} when needed.
\end{itemize}

\hh{Key Takeaways}
\begin{tcolorbox}[% passing options here
    colframe=orange,
    colback = white,
%    colback=orange!40!yellow!30, % think of mixing 2 colors with sliders
    arc=2mm % sharp corners
 ]
\begin{itemize}
\item NumPy arrays enable fast, clear math for geometric grids, loads, and performance studies.
\item Learn to read shape, ndim, and dtype; master slicing and masks.
\item Broadcasting and ufuncs let you express complex calculations in a few lines.
\item Always watch for views—use \c{.copy()} when independence matters.
\end{itemize}
\end{tcolorbox}